\documentclass{ximera}




% MATH -----------------------------------------------------------
\newcommand{\nats}{\mathbb N}
\newcommand{\ints}{\mathbb Z}
\newcommand{\rats}{\mathbb Q}
\newcommand{\reals}{\mathbb R}
\newcommand{\complex}{\mathbb C}
\newcommand{\powerset}{\mathscr P}

%-------------------------------------------------------------

\pgfplotsset{compat=1.18}
\begin{document}
\begin{abstract}
 \end{abstract}

    \title{Class Discussion Activity 21}
    
    \maketitle

\section*{Exercises}

\begin{enumerate}[label=(\arabic*)]


\item Which is the inverse Laplace transform of $\displaystyle F(s)=\frac{s-1}{s^2+2s+5}$?

\begin{enumerate}[label=(\alph*)]
\item $f(t)=\frac{1}{2}e^{-t}\left(2\cos{(2t)}-\sin{(2t)}\right)$
\item $f(t)=e^{-t}\left(\cos{(2t)}-2\sin{(2t)}\right)$
\item $f(t)=\frac{1}{2}e^{-t}\left(2\cos{(2t)}+\sin{(2t)}\right)$
\item $f(t)=e^{-t}\left(\cos{(2t)}+\sin{(2t)}\right)$ 
\item $f(t)=e^{-t}\left(\cos{(2t)}-\sin{(2t)}\right)$  % ***
\item $f(t)=e^{-t}\sin{(2t)}$
\end{enumerate}

% (e) with options through (f)


\item Evaluate $\displaystyle \mathcal{L}^{-1}\left(\frac{s+9}{(s^2+2s+2)(s^2-1)}\right)$ at $t=4$ to the nearest tenth.

% (s+9)/([s^2+2s+2](s-1)(s+1)]  = 1/(s-1)-4/(s+1)+(As+B)/(s^2+2s+2) ---> A=3, B=1.

% So 1/(s-1)-4/(s+1)+(3s+1)/(s^2+2s+2)= 1/(s-1)-4/(s+1)+[3(s+1)-2]/(s^2+2s+2) has inv. transform of

% e^t-4e^(-t)+e^(-t)[3cos(t)-2sin(t)] which at t=4 is approx. 54.5


\item Evaluate $\displaystyle \mathcal{L}^{-1}\left(\dfrac{28s^2+2}{16s^4+s^2}\right)$ at $t=2\pi$.  Round to the nearest thousandth.

% f(t) = 2t -  sin(t/4)

% f(2pi) = 4pi - 1 = 11.566


\item Evaluate $\displaystyle \mathcal{L}^{-1}\left(\frac{s^4-3s^3+18s^2+81}{s^6+18s^4+81s^2}\right)$ at $t=1$ to the nearest hundredth.

    % Ans: f(t) = t-\frac{1}{2}t\sin{(3t)} ---> Laplace: 1/s^2 +0.5 d/(ds) (3/(s^2+9))= 1/s^2 - 3s/(s^2+9)^2

  %  f(1)=1-0.5sin(3)\approx 0.93

\newpage


\item Let $a,b$ be constants, $a>0$, and $f(t)$ be a continuous function. Given $\displaystyle F(s)=\mathcal{L}(f(t))=\int_0^{\infty} f(t)e^{-st}\,dt$, which is $\mathcal{L}^{-1}(F(as-b))$ in terms of $a$, $b$, and the function $f$?

Hint: use $u$-substitution to determine an expression for $F(as)$; then use the First-Shifting Property to determine an expression for\\ $F(a(s-b/a))=F(as-b)$.

% 1/a f(t/a)e^{(b/a)t}


\begin{enumerate}[label=(\alph*)]
\item $\displaystyle f\left(\frac{t}{a}\right)$
\item $\displaystyle f\left(\frac{t}{a}\right)e^{bt}$
\item $\displaystyle f\left(\frac{t}{a}\right)e^{bt/a}$
\item $\displaystyle f\left(\frac{t}{a}\right)e^{-bt}$
\item $\displaystyle f\left(\frac{t}{a}\right)e^{-bt/a}$
\item $\displaystyle \frac{1}{a}f\left(\frac{t}{a}\right)$
\item $\displaystyle \frac{1}{a}f\left(\frac{t}{a}\right)e^{bt}$
\item $\displaystyle \frac{1}{a}f\left(\frac{t}{a}\right)e^{bt/a}$ % ***
\item $\displaystyle \frac{1}{a}f\left(\frac{t}{a}\right)e^{-bt}$
\item $\displaystyle \frac{1}{a}f\left(\frac{t}{a}\right)e^{-bt/a}$
\end{enumerate}


% (h) with options through (j)



%  F(s)=L(f(t))= Int_0^infty f(t)e^(-st)dt ---> F(as)=Int_0^infty f(t)e^(-ast)dt.  Use U-sub.   u=at.  du /a=dt.  F(as) =1/a Int_0^infty f(u/a)e^(-su)du =1/a L(f(t/a))

%    First shifting property: If F(as)=1/a L(f(t/a)) then  L(1/a f(t/a) e^(b/a t))=F(a(s-b/a))=F(as-b).  So the L^(-1) (F(as-b))=1/a f(t/a)e^{(b/a)t}.



\end{enumerate}



\end{document}